\begin{abstract}
Dynamic scene reconstruction in autonomous driving remains a fundamental challenge due to significant temporal variations, moving objects, and complex scene dynamics. 
Existing feed-forward 3D models have demonstrated strong performance in static reconstruction but still struggle to capture dynamic motion. 
To address these limitations, we propose DynamicVGGT, a unified feed-forward framework that extends VGGT from static 3D perception to dynamic 4D reconstruction. 
Our goal is to model point motion within feed-forward 3D models in a dynamic and temporally coherent manner. 
To this end, we jointly predict the current and future point maps within a shared reference coordinate system, allowing the model to implicitly learn dynamic point representations through temporal correspondence. 
To efficiently capture temporal dependencies, we introduce a Motion-aware Temporal Attention (MTA) module that learns motion continuity. 
Furthermore, we design a Dynamic 3D Gaussian Splatting Head that explicitly models point motion by predicting Gaussian velocities using learnable motion tokens under scene flow supervision. It refines dynamic geometry through continuous 3D Gaussian optimization. 
Extensive experiments on autonomous driving datasets demonstrate that DynamicVGGT significantly outperforms existing methods in reconstruction accuracy, achieving robust feed-forward 4D dynamic scene reconstruction under complex driving scenarios.

% We will release the code and pretrained models upon acceptance.
\end{abstract}

% hzl - 1022edit
% v2 从 参考 page-4d改
% 在自动驾驶中，动态场景重建仍然是一项基础且具有挑战性的任务，这主要源于时间维度上的剧烈变化、移动物体的存在以及动态场景对静态几何假设的违背。现有的前馈式三维模型（如 Visual Geometry Grounded Transformer，VGGT）在静态场景重建中表现优异，但在处理时序一致性和动态运动方面存在不足。为此，我们提出了 DynamicVGGT，一个将 VGGT 从静态三维感知扩展至动态四维重建的统一前馈框架。DynamicVGGT 的核心在于一种 动态点图（Dynamic Point Map）表示，该机制在共享参考坐标系下联合预测当前帧与未来帧的三维点图，其位移可用于粗略估计场景流并建模动态几何变化。为了高效捕获时序依赖关系，我们进一步引入了 并行时间注意力（Parallel Temporal Attention） 模块，以学习跨帧的长程运动连续性；同时设计了基于高斯的 场景变形与点追踪（GSDPT） 模块，通过连续三维高斯优化对动态几何进行精细重建。大量在合成与真实自动驾驶数据集上的实验表明，DynamicVGGT 在重建精度与时序一致性方面均取得显著提升，能够在复杂驾驶场景中实现鲁棒的前馈式四维动态场景重建。